% =====================================================================
% Clase -- Física 9° -- Trimestre I -- Semana 04
% Sesión 003 (lun 21 sep., 10:10, 50 min)
% Tema: Científicos que cambiaron la física
% Subtema: Exposiciones de científicos (3 de 7); cierre: describir no es explicar
% Hilo conductor del trimestre I: «De Aristóteles a Galileo: aprender a describir el movimiento»
% Tema 1 de la guía: Científicos que cambiaron la física (tema:cientificos)
% Segundo y último día de exposiciones: en la sesión 002 se presentaron 4 y hoy van las 3
% restantes (decisión de la docente, 2026-09-19). Por eso la sesión no lleva taller.
% Material del docente (no se entrega a los estudiantes).
% =====================================================================
\documentclass[11pt]{article}
\makeatletter\def\input@path{{../../../../../plantillas/estilo/}}\makeatother
\usepackage{mmcantillo}
\guiadelestudiante{../../guia-didactica/guia-periodo-I-movimiento}

\tipodocumento{Clase}
\asignatura{Física}
\titulo{Semana 4: de explicar a medir}
\grado{9°}
\periodo{I}

% DBA: naturales grado 9 · DBA 1 — Comprende que el movimiento de un cuerpo, en un marco de
% referencia inercial dado, se puede describir con gráficos y predecir por medio de
% expresiones matemáticas. Evidencia de esta sesión: sienta la distinción entre describir y
% explicar el movimiento, de la que dependen las tres evidencias del DBA. La columna «me
% aproximo al conocimiento como científico-a natural» de los estándares sostiene la
% exposición y la coevaluación.

\begin{document}

\encabezado*

\begin{center}
{\renewcommand{\arraystretch}{1.3}
\begin{tabularx}{\linewidth}{@{}l l l X l@{}}
  \toprule
  \textbf{Sesión} & \textbf{Fecha} & \textbf{Min} & \textbf{Subtema} & \textbf{Entregables}\\
  \midrule
  003 & lun 21 sep., 10:10 & 50 & Exposiciones (3 de 7) y cierre: describir no es explicar
      & Se recoge la tarea de la sesión 002\\
  \bottomrule
\end{tabularx}}
\end{center}

Cierre del primer tema. En la sesión 002 se presentaron cuatro exposiciones; hoy van las tres
que faltan y, con ellas en la mano, se nombra el cambio que sostiene todo el trimestre: pasar
de \emph{explicar} por qué se mueven los cuerpos, como Aristóteles, a \emph{describir} cómo se
mueven, como Galileo. Referencia: \refguia{tema:cientificos}.

% ---------------------------------------------------------------------
\seccion{Sesión 003 — lunes 21 de septiembre · 10:10 · 50 min}

\textbf{Objetivo.} Presentar las tres exposiciones pendientes y distinguir, con ejemplos de
las siete, entre explicar un fenómeno y describirlo con mediciones.

\textbf{Materiales.} Tablero; el cartel o las diapositivas de cada grupo; la guía; cronómetro
(el del celular sirve) para cuidar los tiempos.

\begin{center}
{\renewcommand{\arraystretch}{1.3}
\begin{tabularx}{\linewidth}{@{}l l X@{}}
  \toprule
  \textbf{Momento} & \textbf{Min} & \textbf{Actividad}\\
  \midrule
  Apertura & 0--4 & Recoger la tarea. Recordar en una frase las cuatro exposiciones de la
    semana pasada y anunciar que hoy se cierra el tema.\\
  Exposiciones & 4--40 & Tres exposiciones de unos 8 minutos, con 3--4 de preguntas cada
    una. Usar la rejilla de abajo mientras exponen; el resto del curso llena la fila que le
    corresponde en la coevaluación.\\
  Puesta en común & 40--47 & La pregunta del día (ver abajo): separar en el tablero, en dos
    columnas, lo que cada científico \emph{explicó} y lo que \emph{midió}.\\
  Cierre & 47--50 & Leer la definición de movimiento uniforme de la guía y anunciar la
    sesión 004: empezamos a describir el movimiento con posición y trayectoria.\\
  \bottomrule
\end{tabularx}}
\end{center}

\textbf{La pregunta del día.} \emph{¿Cuál de estos científicos explicó y cuál midió?}

Se escriben en el tablero dos columnas, «Explicar» y «Describir/medir», y el curso va
ubicando a los siete. La idea no es que la respuesta sea única, sino que se note el cambio de
oficio: Aristóteles da una \emph{causa} (los cuerpos caen porque buscan su lugar natural);
Galileo da una \emph{regla con números} (distancias iguales en tiempos iguales). Dos ayudas
por si el grupo se atasca:

\begin{itemize}
  \item «¿Este señor podría estar equivocado y aun así sonar convincente?» — así se llega a
        que una explicación sin medida no se puede comprobar.
  \item «¿Qué instrumento necesitó?» — el reloj de agua y el plano inclinado de Galileo
        frente a la ausencia de instrumentos en Aristóteles.
\end{itemize}

\textbf{Rejilla para calificar la exposición (5 puntos).}

\begin{center}
{\renewcommand{\arraystretch}{1.3}
\begin{tabularx}{\linewidth}{@{}X l@{}}
  \toprule
  \textbf{Criterio} & \textbf{Puntos}\\
  \midrule
  Dice qué preguntó su científico y qué respondió (la idea, no la biografía) & 1,5\\
  Cuenta al menos un experimento o una observación concreta, con lo que midió & 1,5\\
  Explica con sus palabras, sin leer el cartel & 1,0\\
  Responde las preguntas del curso & 0,5\\
  Se ajusta al tiempo y cita de dónde sacó la información & 0,5\\
  \midrule
  \textbf{Total} & \textbf{5,0}\\
  \bottomrule
\end{tabularx}}
\end{center}

\begin{nota}
\textbf{Si sobra tiempo} (las exposiciones suelen correrse): la aplicación del relámpago y el
trueno de la guía se lee en dos minutos y deja servida la sesión 004 — el sonido como un
movimiento uniforme que sirve de instrumento de medida.

\textbf{Si falta tiempo:} la puesta en común se puede recortar a las dos columnas del tablero
sin discusión, y la pregunta del día pasa al inicio de la sesión 004 como enganche.
\end{nota}

\begin{nota}
\textbf{Para el cierre, si viene al caso.} Los experimentos que la docente hizo en la primera
semana (primera ley y presión atmosférica) son un buen ejemplo de la diferencia: todos vieron
\emph{qué} pasó, pero nadie midió nada. Sirve para preguntar «¿qué habríamos tenido que medir
para convertir eso en física?».
\end{nota}

\begin{nota}
Esta sesión no lleva taller ni tarea nueva: son 50 minutos y las exposiciones ocupan casi
todo. El plan del trimestre pone los dos talleres en las sesiones 005 y 007, y la próxima
tarea en la 004.
\end{nota}

\end{document}
